\documentclass[english, parskip=full]{scrartcl}
\usepackage[utf8]{inputenc}  % Kodierung der Datei
\usepackage[english]{babel}  % Sprache auf Deutsch 
\usepackage{color}
\usepackage{graphicx, gensymb}
\usepackage{amsfonts, cancel, amssymb}
\usepackage{amsmath, amsthm, array, esint}
\usepackage{booktabs, multirow}  % schönere Tabellen
\usepackage{caption}  % Anpassung der Captions
\usepackage{scrlayer-scrpage}    % Manipulation des Seitenstils
\pagestyle{scrheadings}  % Stil für die Seite setzen
\automark{section}  % Abschnittsnamen als Seitenbeschriftung 
\setheadsepline{.5pt}    % Kopzeile durch Linie abtrennen
\usepackage[hidelinks]{hyperref}  % Links und weitere PDF-Features
\usepackage{typearea, tikz, pgffor, verbatim}
\usetikzlibrary{calc, quotes}
\usepackage[cache=false, outputdir=build]{minted}
\usepackage{placeins} % provides \FloatBarrier

\definecolor{bg}{rgb}{0.9,0.9,0.9}
\newcommand\numberthis{\addtocounter{equation}{1}\tag{\theequation}}
\newcommand{\bra}{}
\newcommand{\kom}{}
\newcommand{\brak}{}
\newcommand{\braket}{}
\newcommand{\intx}{}
\newcommand{\intr}{}
\newcommand{\kla}{}
\newcommand{\klam}{}
\newcommand{\klammer}{}
\newcommand{\oper}{}
\newcommand{\tecxt}{}
\newcommand{\Bra}{}
\newcommand{\Ket}{}
\newcommand{\I}{\text{i}}
\renewcommand{\i}{\text{i}}
\newcommand{\e}{\text{e}}
\newcommand*{\Fourier}{\textsc{Fourier}\ }
\newcommand*{\F}{\mathcal{F}}
\newcommand*{\sinc}{\text{sinc\,}}
\newcommand{\conv}{\mathop{\scalebox{3.5}{\raisebox{-0.38ex}{\makebox[0.5\width][c]{$\ast$}}}}\limits}

\newcommand*{\titel}{Variants on Formulae for Frequency Analysis}
\newcommand*{\autor}{Joachim Schwardt}

\hypersetup{pdfauthor={\autor}, pdftitle={\titel}}

%\titlehead{titlehead}
%\subject{subject}
\title{\titel}
%\author{\autor}
\date{\today}

\begin{document}

%\begin{titlepage}
\maketitle  % Titel setzen
%\tableofcontents  % Inhaltsverzeichnis setzen
%\end{titlepage}

\section{Discrete Fourier Transformation with Hann-Window}
NOTE: The following is copied from the original file 'Naff{\_}formulae.tex'.
\\
Model for signal:
\begin{align}
  z^{\text{model}}_t &= \sum_k c_k e^{i 2 \pi \nu k t}
  \qquad \text{ with e.g. }
  k \in \{-3,-2,-1,0,1,2,3\}
  \label{eq_model}
\end{align}
Hann-Window:
\begin{align}
  w_t &= 2 \sin^2 \frac{\pi t}{N} = 1 -\cos \frac{2\pi t}{N}
       = 1 - \frac{1}{2} e^{i 2 \pi \frac{t}{N}} - \frac{1}{2} e^{- i 2 \pi \frac{t}{N}}
\end{align}
Will be used with $t=0,1,\dots,N-1$ (as it keeps the formulae simpler)
instead of $t=1,\dots,N$ (as in [2]).

Discrete Fourier transform (assumes $\nu$ irrational in Eq.(\ref{eq_geomseries})):
\begin{align}
  F^{\text{model}}_j &= \frac{1}{N} \sum_{t=0}^{N-1} z^{\text{model}}_t w_t e^{- i 2 \pi \frac{j t}{N}}
  \qquad
  j=0,1,\dots, N-1
  \\
  &=
  \frac{1}{N} \sum_{t=0}^{N-1} \sum_k c_k e^{i 2 \pi \nu k t}
  \left( 1 - \frac{1}{2} e^{i 2 \pi \frac{t}{N}} - \frac{1}{2} e^{- i 2 \pi \frac{t}{N}} \right)
  e^{- i 2 \pi \frac{j t}{N}}
  \\
  &=
  \frac{1}{N} \sum_k c_k \sum_{t=0}^{N-1} \left[
                  \left( e^{i 2 \pi (\nu k - \frac{j}{N}) } \right)^t
    - \frac{1}{2} \left( e^{i 2 \pi (\nu k - \frac{j-1}{N}) } \right)^t
    - \frac{1}{2} \left( e^{i 2 \pi (\nu k - \frac{j+1}{N}) } \right)^t
  \right]
  \\
  &=
  c_0 (\delta_{j,0} - \frac{1}{2} \delta_{j,1} - \frac{1}{2} \delta_{j,N-1} )
  \\
  &
  + \frac{1}{N} \sum_{k \ne 0} c_k \left[
                  \frac{1-e^{i 2 \pi (\nu k - \frac{j}{N})N }}{1-e^{i 2 \pi (\nu k - \frac{j}{N}) } }
    - \frac{1}{2} \frac{1-e^{i 2 \pi (\nu k - \frac{j-1}{N})N }}{1-e^{i 2 \pi (\nu k - \frac{j-1}{N}) } }
    - \frac{1}{2} \frac{1-e^{i 2 \pi (\nu k - \frac{j+1}{N})N }}{1-e^{i 2 \pi (\nu k - \frac{j+1}{N}) } }
  \right]
  \label{eq_geomseries}
  \\
  &=
  c_0 (\delta_{j,0} - \frac{1}{2} \delta_{j,1} - \frac{1}{2} \delta_{j,N-1} )
  \\
  &
  + \frac{1}{N} \sum_{k \ne 0} c_k (1-e^{i 2 \pi \nu k N })\left[
                  \frac{1}{1-e^{i 2 \pi (\nu k - \frac{j}{N}) } }
    - \frac{1}{2} \frac{1}{1-e^{i 2 \pi (\nu k - \frac{j-1}{N}) } }
    - \frac{1}{2} \frac{1}{1-e^{i 2 \pi (\nu k - \frac{j+1}{N}) } }
  \right]
  \\
  &= \dots
\end{align}
%\\Note: checked by Ketzmerick and Lange (FIXME: also check validity for resonant $\nu$ of Eqs.~(\ref{eq_dfft},\ref{r_frac}))
Final result
\begin{equation}
  \fbox{$
  F^{\text{model}}_j =
  c_0 (\delta_{j,0} - \frac{1}{2} \delta_{j,1} - \frac{1}{2} \delta_{j,N-1} )
  - \frac{\sin^2 \frac{\pi}{N}}{N} {\displaystyle \sum_{k \ne 0}^{}} c_k
  e^{i \pi \nu k N} \sin(\pi \nu k N) \;\;
  \frac{\cos[\pi(\nu k - \frac{j}{N})]}
       {\sin[\pi(\nu k - \frac{j}{N})] \sin[\pi(\nu k - \frac{j+1}{N})] \sin[\pi(\nu k - \frac{j-1}{N})]}
  $}
  \label{eq_dfft}
\end{equation}



\clearpage
\section{Compute Frequencies}
\subsection{Recalculation of Analytic Formula from [2]}
Summary: The result of [2] is extended to frequencies close to zero, $\nu < \frac 1N$,
or close to one, $\nu > 1-\frac 1N$.

Note: The shifted use of the Hann-Window compared to [2] is
of no relevance for the result.

Use just one oscillatory component $c_1 \ne 0$ in model:
\begin{eqnarray}
  z^{\text{model}}_t &=& c_0 + c_1 e^{i 2 \pi \nu t}
\end{eqnarray}
For a given signal $z_t$ the frequency $\nu$ has to be determined with $c_0$ and $c_1$ unknown.
Use from Eq.~(\ref{eq_dfft})
\begin{align}
  F^{\text{model}}_j &=
  c_0 (\delta_{j,0} - \frac{1}{2} \delta_{j,1} - \frac{1}{2} \delta_{j,N-1} )
  - \\
  &-c_1 \frac{\sin^2 \frac{\pi}{N}}{N}
  e^{i \pi \nu N} \sin(\pi \nu N) \;\;
  \frac{\cos[\pi(\nu - \frac{j}{N})]}
       {\sin[\pi(\nu - \frac{j}{N})] \sin[\pi(\nu - \frac{j+1}{N})] \sin[\pi(\nu - \frac{j-1}{N})]}
\end{align}
Take the negative ratio of any two neighboring values at $j,j+1$ which are not from $\{0,1,N-1\}$, as $c_0$ is not known.
\begin{eqnarray}
  \label{r_frac}
  R &:=&
  - \frac{F^{\text{model}}_j}{F^{\text{model}}_{j+1}} =
  - \frac{\cos[\pi(\nu - \frac{j}{N})]   \sin[\pi(\nu - \frac{j+2}{N})]}
       {\cos[\pi(\nu - \frac{j+1}{N})] \sin[\pi(\nu - \frac{j-1}{N})]}
  \label{eq_defR}
\end{eqnarray}
This value is thus real without taking the absolute value of the Fourier transform (as in [2]).
Its sign is positive if $\frac{j-1}{N} < \nu < \frac{j+2}{N}$ otherwise negative
(assuming $\nu$ to be somewhere around $\frac{j}{N}$ such that $\cos[\pi(\nu - \frac{j}{N})]$
is still positive).

Define for brevity
\begin{eqnarray}
  x &:=& \pi(\nu - \frac{j}{N})
  \qquad
  b:= \frac{\pi}{N}
\end{eqnarray}
such that we have (using $\sin a \cos b = \frac12 [\sin(a-b) + \sin(a+b)]$)
\begin{eqnarray}
  R &=&
  - \frac{\cos[x] \sin[x-2b]}{\cos[x-b] \sin[x+b]}
  =
  - \frac{-\sin[2b] + \sin[2x-2b]}{\sin[2b] + \sin[2x]}
  \\
  &=&
  \frac{-\sin 2x \cos 2b  + (1+\cos 2x) \sin 2b }
       {\sin 2x  + \sin 2b}
  \\
  &=&
  \frac{-\frac{\sin 2x}{\sin 2b} \cos 2b  + 1+\cos 2x}
       {\frac{\sin 2x}{\sin 2b}+1}
  =
  \frac{-c y  + 1+\cos 2x}
       {y+1}
\end{eqnarray}
where we defined
\begin{eqnarray}
  \label{y_def}
  y &:=& \frac{\sin 2x}{\sin 2b} = \frac{\sin 2\pi(\nu - \frac{j}{N})}{\sin 2b}
  \qquad
  c:=\cos 2b = \cos \frac{2\pi}{N}
\end{eqnarray}
This gives
\begin{eqnarray}
  R &=&
  \frac{-c y  + 1 + \sqrt{1-(1-c^2)y^2}}{y  + 1}
  \label{eq_R}
\end{eqnarray}
and after rearrangement
\begin{eqnarray}
  y(R+c) + (R-1) &=&
  \sqrt{1-(1-c^2)y^2}
  \label{eq_Rsq}
\end{eqnarray}
leading after squaring to a quadratic equation in $y$ with two solutions
\begin{eqnarray}
  \label{eq_yofR}
  y_\pm &=&
  \frac{-(R+c)(R-1) \pm \sqrt{c^2(R+1)^2 - 2R(2c^2-c-1)}}
       {R^2 + 1 + 2Rc}
\end{eqnarray}
This corresponds to Eq.(26) in [2] with $R=a/b$ and the positve sign before the square root.

May it happen that the solutions are complex? No, this would happen for $R \approx -1$ only,
which however is not possible according to Eq.~(\ref{eq_defR}) for $\frac jN$ close to $\nu$.

Which of the two solutions has to be considered?
The solution which gives a positive l.h.s.\ in
Eq.~(\ref{eq_Rsq}) corresponding to the positively defined square root on the r.h.s., i.e.
\begin{eqnarray}
  \label{y_plus_minus}
  0 &\leq& y_\pm (R+c) + (R-1)
\end{eqnarray}
This inequality can be evaluated by multiplying with the positive factor
$R^2 + 1 + 2Rc = (R+c)^2 + s^2$
\begin{align}
  0 &\leq [y_\pm (R+c) + (R-1)] (R^2 + 1 + 2Rc)
  =\\
  &= \pm (R+c) \sqrt{c^2(R+1)^2 - 2R(2c^2-c-1)} + s^2(R-1)
\end{align}
Again using that $R \approx -c \approx -1$ does not happen and that the second term is small,
the sign of $(R+c)$ decides about the correct solution,
\begin{align}
  y &= \left\{
          \begin{array}{l}
             y_+ \quad \text{for } R+c > 0
             \\
             y_- \quad \text{otherwise}
          \end{array}
    \right\}
    =\\
    &=
  \frac{-(R+c)(R-1) + \tecxt\text{sign}(R+c) \sqrt{c^2(R+1)^2 - 2R(2c^2-c-1)}}
       {R^2 + 1 + 2Rc}
  \label{eq_ypm_decide}
\end{align}

Note: $R+c>0$ is the case if $\nu > \frac{j-1}{N}$, according to Eq.~(\ref{eq_defR})
and $R+c<0$ is the case if $\nu < \frac{j-1}{N}$ (which applies if $\nu < \frac 1N$).

We now get the frequency with $y$ from Eqs.~(\ref{y_def}) and
(\ref{eq_ypm_decide})
\begin{eqnarray}
  \nu &=&  \frac{j}{N} + \frac{1}{2\pi}\arcsin\left(y \, \sin 2b \right)
\end{eqnarray}

The decision according to Eq.~(\ref{eq_ypm_decide})
works perfectly for a signal of the type considered even for small $\nu < \frac{1}{N}$.
\section{Integral approximation}
NOTE: The following is the 'new' part.
\\
Analysis of the WBA shows, that the specific structure of the used weights may strongly influence the resulting convergence behavior.
For example, using gaussian weights ($\alpha\gtrsim 140$ for $w_{0,N-1}\simeq 10^{-16}$)
\begin{align}
w^\tecxt\text{gauss}_t&=\frac{1}{C}\e^{-\alpha \kla\left( {\frac{t}{N}-\frac{1}{2}} \right)^2}
\end{align}
with the normalization constant
\begin{align}
C&=\sum_{t=0}^{N-1}\e^{-\alpha \kla\left( {\frac{t}{N}-\frac{1}{2}} \right)^2}\simeq N\sqrt{\frac{\pi}{\alpha}}
\end{align}
usually gives better results than the $C_c^\infty$-weights defined by
\begin{align}
w^{C_c^\infty}_t&=\frac{1}{\tilde{C}}\e^{-\klam\left[ {\frac{t}{N}\kla\left( {1-\frac{t}{N}} \right)} \right]^{-1}}.
\end{align}
This gives rise to the idea of also using gaussian weights for the Naff.
While an explicit solution of the resulting summation is not possible for these weights, an integral approximation can be used.
\subsection{Hanning weights}
Consider in the following a slightly simpler model than above, i.e. only the coefficient $c_1$ is to be non-zero.
We repeat the calculation of the Fourier coefficients with the Hanning-weights using this approximation.
This allows us to compare the result of the 'exact Hanning-Naff' (EHN) to the 'approximated Hanning-Naff' (AHN).
It will also demonstrate the approach used for the gaussian weights.
\\
The coefficients can be calculated by
\begin{align}
F_j&=\sum_{t=0}^{N-1}x_t^\tecxt\text{model}w_t\e^{-2\pi\I j\frac{t}{N}}=\\
&=\frac{c_1}{N}\sum_{t=0}^{N-1}\kla\left( {1-\frac{1}{2}\e^{2\pi\I\frac{t}{N}}-\frac{1}{2}\e^{-2\pi\I \frac{t}{N}}} \right)\e^{2\pi\I \frac{t}{N}\kla\left( {\nu N - j} \right)}=\\
&=\frac{c_1}{N}\sum_{l\in\{0,-1,+1\},\,\beta_l\in\{1,-\frac{1}{2},-\frac{1}{2}\}}\sum_{t=0}^{N-1}\e^{2\pi\I \frac{t}{N}\kla\left( {\nu N - j+l} \right)}\approx\\
&\approx c_1\sum_{l\in\{0,-1,+1\},\,\beta_l\in\{1,-\frac{1}{2},-\frac{1}{2}\}}\beta_l\intx\int_{0}^{1}\e^{2\pi\I x (\nu N - j+\beta_l)}\,\mathrm{d}x=\\
&=c_1\sum_{l\in\{0,-1,+1\},\,\beta_l\in\{1,-\frac{1}{2},-\frac{1}{2}\}}\beta_l\frac{\e^{2\pi\I (\nu N-j+l)}-1}{2\pi\I (\nu N-j+l)}=\\
&=\frac{c_1}{2\pi\I}\kla\left( {\e^{2\pi\I \nu N}-1} \right)\sum_{l\in\{0,-1,+1\},\,\beta_l\in\{1,-\frac{1}{2},-\frac{1}{2}\}}\klam\left[ {\frac{1}{\nu N-j}-\frac{1}{2}\frac{1}{\nu N -j+1}-\frac{1}{2}\frac{1}{\nu N-j-1}} \right].
\end{align}
Note that 
\begin{align}
\frac{1}{a}-\frac{1}{2}\frac{1}{a+1}-\frac{1}{2}\frac{1}{a-1}&=\frac{1}{a}-\frac{a}{(a+1)(a-1)}=-\frac{1}{a(a+1)(a-1)}.
\end{align}
Inserting this into the above with $a=\nu kN-j$ gives
\begin{align}
F_j&=\frac{c_1}{2\pi\I}\kla\left( {\e^{2\pi\I \nu N}-1} \right)\frac{-1}{(\nu N-j)(\nu N-j+1)(\nu N-j-1)}.
\end{align}
We are again interested in the ratio of consecutive Fourier coefficients
\begin{align}
R&=\frac{F_j}{F_{j+1}}\approx\frac{(\nu N-(j+1))(\nu N-(j+1)+1)(\nu N-(j+1)-1)}{(\nu N-j)(\nu N-j+1)(\nu N-j-1)}=\frac{\nu N-j-2}{\nu N -j +1}.
\end{align}
Solving for $\nu$ leads to
\begin{align}
R(\nu N -j +1)&=\nu N-j-2\\
\nu N (R-1)&=R(j-1)-j-2\\
\nu&=\frac{j}{N}-\frac{1}{N}\frac{2+R}{R-1}. \label{eqn:ApproxHanningNaff}
\end{align}
Figure \ref{figure:Naff-hann-comp} shows that the approximation is, from a practical standpoint, completely equivalent to the exact solution.
In fact, due to a certain ambiguity about $\nu$ vs. $1-\nu$, the more simple approximation may even be for choice in certain cases.
\begin{figure}[h!]
\centering
\includegraphics[width=\linewidth]{pictures/Naff_hann_comp.png}
\caption{Comparison of the exact Hanning Naff (EHN) and the approximated one (AHN).
In each case, the convergence behavior to the frequency computed using the respective formula is shown.
Note that a correction is needed for the EHN, since sometimes the detected frequency is not $\nu$ but rather $1-\nu$.
Interestingly, this is not required at all for the AHN.}
\label{figure:Naff-hann-comp}
\end{figure}
If anything, this result gives hope that a similar solution can be found for gaussian weights, for which an exact and explicit solution is not attainable.
\newpage
\subsection{Gaussian weights}
Transforming the relevant sum over $t$ into an integral over $x\in[0,1]$ will only be of any help, if we can approximate the integral by expanding this interval to $\mathbb{R}$.
This approximation requires $\alpha$ to be large enough, such that this error vanishes within the machine precision. 
%(TODO: check if smaller $\alpha$ really give worse results. Since we are mostly just interested in the ratio $R$, errors may cancel out.)
\\
We start by again computing the Fourier coefficients
\begin{align}
F_j&=\sum_{t=0}^{N-1}x_t^\tecxt\text{model}w_t\e^{-2\pi\I j\frac{t}{N}}=\\
&=\frac{c_1}{C}\sum_{t=0}^{N-1}\e^{-\alpha\kla\left( {\frac{t}{N}-\frac{1}{2}} \right)^2}\e^{2\pi\I \frac{t}{N}\kla\left( {\nu N - j} \right)}\approx\\
&\approx\frac{c_1N}{C}\intx\int_{0}^{1}\e^{-\alpha\kla\left( {x-\frac{1}{2}} \right)^2}\e^{2\pi\I x\kla\left( {\nu N - j} \right)}\,\mathrm{d}x=\\
&=\frac{c_1N}{C}\e^{\pi\I (\nu N-j)}\intx\int_{-1/2}^{1/2}\e^{-\alpha x^2}\e^{2\pi\I x\kla\left( {\nu N - j} \right)}\,\mathrm{d}x\approx\\
&\approx\frac{c_1N}{C}\e^{\pi\I (\nu N-j)}\intr\int_{\mathbb{R}^{ }}\e^{-\alpha x^2}\e^{2\pi\I x\kla\left( {\nu N - j} \right)}\,\mathrm{d}^{ }x\approx\\
&\approx c_1\e^{\pi\I (\nu N-j)}\e^{\frac{\left[2\pi\I(\nu N-j)\right]^2}{4\alpha}}=c_1\e^{\pi\I (\nu N-j)}\e^{-\frac{\pi^2}{\alpha}(\nu N-j)^2}. \label{eqn:Fj-gauss}
\end{align}
The ratio of consecutive coefficients then becomes
\begin{align}
R&=\left\vert\frac{F_j}{F_{j+1}}\right\vert=\e^{\frac{\pi^2}{\alpha}\klam\left[ {(\nu N-(j+1))^2-(\nu N-j)^2} \right]}.
\end{align}
With $a=\nu N-j$ and $(a-1)^2-a^2=1-2a$ this can be reduced to
\begin{align}
\frac{\alpha}{\pi^2}\log R&=1-2(\nu N-j)\\
\nu&=\frac{j}{N}+\frac{1}{2N}-\frac{\alpha}{2\pi^2N}\log R. \label{eqn:ApproxGaussNaff}
\end{align}
This result will also be referred to as the 'approximated Gauss-Naff' (AGN).
\\
More generally, we define the ratios
\begin{align}
R_k &:= \left\vert\frac{F_j}{F_{j+k}} \right\vert = \e^{\frac{\pi^2}{\alpha}\klam\left[ {(\nu N-(j+k))^2-(\nu N-j)^2} \right]}.
\end{align}
With $a=\nu N-j$ and $(a-k)^2-a^2=k^2-2ak$ this can be reduced to
\begin{align}
\frac{\alpha}{\pi^2}\log R_{j,k}&=k^2-2k(\nu N-j)\\
\nu_{j,k}&=\frac{j}{N}+\frac{k}{2N}-\frac{\alpha}{2k\pi^2N}\log R_{j,k}. \label{eqn:ApproxGaussNaff k ratio}
\end{align}
Note that we may introduce a series of frequencies
\begin{align}
\nu'_l &:= \nu_{j-l,2l+1} =\frac{j-l}{N} + \frac{2l+1}{2N}-\frac{\alpha}{2(2l+1)\pi^2N}\log \left\vert \frac{F_{j-l}}{F_{j+l+1}} \right\vert
\end{align}
for integers $l=\{0,1,\dots\}$.
\\
The following figures show a number of examples of the convergence of $\nu_N$ to the true frequency $\nu$.
For most of them, the WBA with $C_c^\infty$-weights and gaussian weights is shown for comparison.
Since the results of the two Hanning-Naff variants are practically identical, only the EHN is shown.
At relevant sections, the linear slope is shown.
Since the plots are logarithmic, this coincides with a power-law decay.
\begin{figure}[h!]
\centering
\includegraphics[width=\linewidth]{pictures/Naff_variants.png}
\caption{Comparison of Hanning-Naff and Gauss-Naff.}
\label{figure:Naff-variants}
\end{figure}
\begin{figure}[h!]
\centering
\includegraphics[width=\linewidth]{pictures/Naff_WBA_variants.png}
\caption{Comparison of Hanning-Naff, Gauss-Naff and both WBAs.}
\label{figure:Naff-WBA-variants}
\end{figure}
\begin{figure}[h!]
\centering
\includegraphics[width=\linewidth]{pictures/Naff_WBA_variants2.png}
\caption{Comparison of Hanning-Naff, Gauss-Naff and both WBAs.}
\label{figure:Naff-WBA-variants2}
\end{figure}
\begin{figure}[h!]
\centering
\includegraphics[width=\linewidth]{pictures/Naff_WBA_variants3.png}
\caption{Comparison of Hanning-Naff, Gauss-Naff and both WBAs.}
\label{figure:Naff-WBA-variants3}
\end{figure}
\begin{figure}[h!]
\centering
\includegraphics[width=\linewidth]{pictures/Naff_WBA_variants4.png}
\caption{Comparison of Hanning-Naff, Gauss-Naff and both WBAs.}
\label{figure:Naff-WBA-variants4}
\end{figure}
\FloatBarrier
\newpage
\subsection{Nonzero offset of the signal}
There remains the issue of assuming $c_0=0$.
In the case of the standard map, the offset of the center island may be known by the used convention of the map.
However, it is also possible to solve this issue for arbitrary offsets.
The same derivation that lead to eqn. (\ref{eqn:Fj-gauss}) can be repeated for the case of $c_0\neq 0$.
The $c_0$-contribution to $F_j$ is therefor given by
\begin{align}
F_j^{(0)}&\approx\frac{c_0N}{C}\intx\int_{0}^{1}\e^{-\alpha\kla\left( {x-\frac{1}{2}} \right)^2}\e^{-2\pi\I xj}\,\mathrm{d}x\approx\\
&\approx\dots\approx c_0\e^{-\pi\I j}\e^{-\frac{\pi^2}{\alpha}j^2}. \label{eqn:Fj-c0-corretion}
\end{align}
All that is needed to apply the previous result (\ref{eqn:ApproxGaussNaff}) is to transform the Fourier coefficients according to
\begin{align}
F_j&\longrightarrow F_j-F_j^{(0)}. \label{eqn:Fj-gauss-transform}
\end{align}
Figure \ref{figure:Naff-offset-comp} shows how this transformation gives the same result as subtracting the 'correct' offset of $c_0=\frac{1}{2}$ for an oscillatory signal of the standard map.
\begin{figure}[h!]
\centering
\includegraphics[width=\linewidth]{pictures/Naff_offset_comp.png}
\caption{Comparison of 'manually' subtracting the correct offset of $c_0=\frac{1}{2}$ from the original signal vs. applying the transformation from equation (\ref{eqn:Fj-gauss-transform}).
Note how the two become practically identical around $N\gtrsim 128$.}
\label{figure:Naff-offset-comp}
\end{figure}
Also note that the value of $c_0$ can be efficiently  approximated by
\begin{align}
c_0&\approx\sqrt{\frac{\alpha}{\pi}}F_0. \label{eqn:c0-gauss-approx}
\end{align}
A similar formula can also be used for the other coefficients $c_k$, corresponding to the $k$-th 'strongest' frequency $\nu_k$ in the signal, leading to
\begin{align}
c_k&\approx\sqrt{\frac{\alpha}{\pi}}F_j\e^{-\I\pi (\nu N-j)}\e^{\frac{\pi^2}{\alpha}\kla\left( {\nu N - j} \right)^2}. \label{eqn:ck-gauss-approx}
\end{align}
Here $j\in \mathbb{N}$ has to be chosen such that $F_j$ is a local maximum of the Fourier coefficients.
\newpage
Given below is the Python code for computing eqn. (\ref{eqn:ApproxGaussNaff}) for an arbitrary signal.
In case the offset of the input signal is known beforehand, the transformation of Fourier-coefficients from (\ref{eqn:Fj-gauss-transform}) is optional and subject to the bool \texttt{Offset}.
Since the correction terms are proportional to $\e^{-\frac{\pi^2}{\alpha}j^2}$, for $\alpha\simeq 140$ already $j\gtrsim 23$ gives a contribution below $10^{-16}$.
Hence $J=30$ as the maximal $j$ is a reasonable cutoff.
\begin{minted}[frame=lines, bgcolor=bg, fontsize=\footnotesize, linenos]{python}
def _remove_peak(fft, ind, nu, N, alpha=140, J=30):
    """Remove a gaussian peak from a FFT at a given frequency"""
    ampl = fft[ind]
    phase = nu * N - ind
    for j in range(-J+1, J, 1):
        corr_exp = -np.pi**2 / alpha * (j**2 - 2 * phase * j)
        fft[ind+j] -= ampl * np.exp(corr_exp)
        
def naffnd_gauss(z, n_freq=1, alpha=140, J=30, 
                 ReturnCoeff=False, Offset=False):
    """
    NaffND with gaussian weights.
    Computes the first 'n_freq' fundamental frequencies of the complex-valued
        input signal 'z', as well as their respective complex amplitudes.
    Parameters:
        n_freq == number of frequenices to be computed until max(fft) < 1e-15
                  (iterative process is aborted if the fft becomes 'flat')
        alpha == gaussian weights given by exp(-alpha * (t / N - 0.5)**2)
                 where t = 0,...,N-1 with 'N' being the signal length
                 (alpha = 140 leads to weights being about 1e-16 at t = 0)
        J == for the '2*J - 1' points closest to the relevant peak in the fft,
             the respective fft-value is subtracted based on an approximate
             analytical formula
             (J = 23 would be enough for machine precision at the outer parts)
        ReturnCoeff == whether to return the computed frequency amplitudes
        Offset == whether to compute, return and adjust a signal offset
    """
    N = z.shape[0]
    t = np.arange(N)    # t = 0,...,N-1
    weights = np.exp(-alpha * (t / N - 0.5)**2)
    fft = np.fft.fft(weights * z)   # need complex-valued fft for coeff.
    abs_fft = np.abs(fft)           # only need absolute values for freqs.
    
    # frequencies and their respective amplitudes
    nu_arr = np.zeros(n_freq + Offset)
    c_arr = np.zeros(n_freq + Offset, dtype=np.complex128)     
    
    if J > N//2:
        J = N//2
    if Offset:    # compute constant offset and adjust fft spectrum
        c_arr[n_freq] = np.sqrt(alpha / np.pi) * fft[0] / N    
        _remove_peak(abs_fft, ind=0, nu=0.0, N=N, J=J)
    
    for i in range(n_freq):
        ind = np.argmax(abs_fft)
        if abs_fft[ind] < 1e-15:   
            return nu_arr   # abort if fft-peak too small -> nu == 0
        if ind == (N-1):   
            ind -= 1    # avoid index error at 'abs_fft[ind + 1]'
        elif abs_fft[ind - 1] > abs_fft[ind + 1] + 1e-8:   
            ind -= 1    # 'ind' closer to true peak (1e-8 tolerance)
            
        # Ratio 'R' of largest neighboring fft-values to compute 'nu'
        R = abs_fft[ind] / abs_fft[ind + 1]
        nu = ind / N + 1 / (2*N) - alpha / (2 * np.pi**2 * N) * np.log(R)
        nu_arr[i] = nu
        
        # compute the complex value of the frequency amplitude
        phase = nu * N - ind
        c = np.sqrt(alpha / np.pi) * fft[ind] / N * np.exp(-1j*np.pi * phase)
        c_arr[i] = c * np.exp(np.pi**2 / alpha * phase**2)
        
        # remove the previous peak from the fft spectrum
        _remove_peak(abs_fft, ind, nu, N, alpha=alpha, J=J)
    if ReturnCoeff:
        return nu_arr, c_arr
    return nu_arr
\end{minted}

\subsection{Flat-top window approximation}
We may define a flat-top window parameterized by
\begin{align}
w(x) &:= \frac{1}{1 + \e^{-\alpha (x-\beta)}} + \frac{1}{1 + \e^{\alpha (x-[1-\beta])}} - 1 = \\
&= \frac{1 + \e^{\alpha (x-[1-\beta])} + 1 + \e^{-\alpha (x-\beta)}}{\kla\left( 1 + \e^{-\alpha (x-\beta)} \right) \kla\left( 1 + \e^{\alpha (x-[1-\beta])} \right)} - 1 = \\
&= \frac{1 - \e^{\alpha (2\beta -1)}}{1 + \e^{-\alpha (x-\beta)}  + \e^{\alpha (x-[1-\beta])} + \e^{\alpha (2\beta -1)}}.
\end{align}
Similarly to the gaussian example, we will again need to assume that $w(x)\approx 10^{-16}$ for $x\in\{0,1\}$, which is coincidentally again satisfied by $\alpha\approx 140$.
The second parameter $\beta$ can be used to adjust, how much of the signal should be ``cut off'', since this window function decays very rapidly for $x < \beta$ and $x > 1 - \beta$ respectively.
We compute the Fourier coefficients using our integral approximation
\begin{align}
F_j&\approx c_1\intx\int_{\mathbb{R}}^{} w(x) \e^{2\pi\I x\kla\left( {\nu N - j} \right)}\,\mathrm{d}x=\\
&=c_1\e^{\pi\I \kla\left( {\nu N - j} \right)} \intr\int_{\mathbb{R}^{ }} \frac{1 - \e^{\alpha (2\beta -1)}}{1 + \e^{\alpha (2\beta -1)} + \e^{\alpha (\beta-1/2)}\kla\left( \e^{\alpha x} + \e^{-\alpha x} \right)} \e^{2\pi\I x\kla\left( {\nu N - j} \right)} \, \mathrm{d}^{ }x = \\
&= c_1\kappa\e^{\pi\I \kla\left( {\nu N - j} \right)} \intr\int_{\mathbb{R}^{ }} \frac{\e^{2\pi\I x\kla\left( {\nu N - j} \right)}}{\gamma +  \e^{\alpha x} + \e^{-\alpha x} }  \, \mathrm{d}^{ }x. \label{eqn:Fj-flattop-integral}
\end{align}
Here we introduced 
\begin{align}
\kappa &:= \e^{\alpha(1/2 - \beta)} - \e^{-\alpha(1/2 - \beta)},\\
\gamma &:= \e^{\alpha(1/2 - \beta)} + \e^{-\alpha(1/2 - \beta)},
\end{align}
where the two parameters could also be expressed using hyperbolic functions.
Let $\phi := \frac{\nu N-j}{\alpha}$ and substitute $x\rightarrow \alpha x$.
Then we are left with the integral
\begin{align}
\alpha I &:= \intr\int_{\mathbb{R}^{ }} \frac{\e^{2\pi\i\phi x}}{\gamma + \e^x + \e^{-x}} \, \mathrm{d}^{ }x
\end{align}
\begin{align}
\intr\int_{\mathbb{R}^{ }} \frac{\e^{\i\omega x}}{1 + \e^{-\alpha(x-\beta)}} \, \mathrm{d}^{ }x &= \e^{\i\omega\beta}\intr\int_{\mathbb{R}^{ }} \frac{\e^{\i\frac{-\omega}{\alpha} y}}{1 + \e^{y}} \, \mathrm{d}^{ }y =  \frac{\i\pi\e^{\i\omega\beta}}{\cosh\kla\left( \frac{\pi\omega}{\alpha} \right)}, \\
\intr\int_{\mathbb{R}^{ }} \frac{\e^{\i\omega x}}{1 + \e^{\alpha (x-[1-\beta])}} \, \mathrm{d}^{ }x &= \frac{\i\pi\e^{\i\omega[1-\beta]}}{\cosh\kla\left( \frac{\pi\omega}{-\alpha} \right)}.
\end{align}
The Fourier transform of the flat-top window function is
\begin{align}
F_j &\approx c_1\frac{\i\pi}{\cosh\kla\left( \frac{2\pi^2[\nu N-j]}{\alpha} \right)}\kla\left( \e^{2\pi\i [\nu N-j]\beta} + \e^{2\pi\i [\nu N-j][1 - \beta]} \right)
\end{align}

%\begin{thebibliography}{9}
%\bibitem{example} description, \url{https://www.exmapleurl}, day.\,month~2021.
%\end{thebibliography}

\end{document}